\documentclass[a4paper,12pt]{article}
\usepackage[ngerman]{babel}
\usepackage[utf8]{inputenc}
\usepackage[T1]{fontenc}
\usepackage{amsmath,amssymb}
\usepackage{tikz}
\usetikzlibrary{automata,positioning,arrows.meta}

\tikzset{
    ->,>=Stealth,
    node distance=2.5cm,
    every state/.style={thick,minimum size=1cm},
    initial text={},
}

\newcommand{\methodbox}[1]{%
\medskip
\noindent\fbox{%
    \begin{minipage}{0.94\linewidth}
    \textbf{Methodik / Definition.}\par\smallskip
    \small
    #1
    \end{minipage}%
}
\medskip\par
}

\newcommand{\examplebox}[1]{%
\medskip
\noindent\fbox{%
    \begin{minipage}{0.94\linewidth}
    \textbf{Beispiele.}\par\smallskip
    \small
    #1
    \end{minipage}%
}
\medskip\par
}

\newcommand{\taskbox}[1]{%
\medskip
\noindent\fbox{%
    \begin{minipage}{0.94\linewidth}
    \textbf{Aufgabenstellung.}\par\smallskip
    \small
    #1
    \end{minipage}%
}
\medskip\par
}

\begin{document}
\raggedright
\setlength{\parindent}{0pt}

{\LARGE\bfseries ETI -- Aufgabenblatt 5 \\ Hausaufgaben\par}
\vspace{2em}

\section*{Aufgabe 5.6}

\taskbox{
Bringen Sie die Grammatik
\[
S\to SaSb\mid A,\qquad
A\to \varepsilon\mid B\mid a,\qquad
B\to bB\mid A\mid C,\qquad
C\to c
\]
in Chomsky-Normalform. Das Startsymbol ist $S$.
}

\methodbox{
Eine CFG ist in Chomsky-Normalform, wenn jede Regel eine der Formen
\[
    A\to BC,\qquad A\to a,\qquad S_0\to\varepsilon
\]
hat. Die letzte Regel ist nur fuer die Startvariable erlaubt, wenn $\varepsilon$ in der Sprache liegt.
Damit die Startvariable nicht rechts vorkommt, fuehrt man zuerst eine neue Startvariable ein.
Danach entfernt man $\varepsilon$-Regeln, Kettenregeln und ersetzt Terminale in laengeren rechten
Seiten durch eigene Variablen.
}

\textbf{Schritt 1: Neue Startvariable.}

Da $S$ rechts in $S\to SaSb$ vorkommt und $\varepsilon$ erzeugt werden kann, fuehren wir eine neue
Startvariable $S_0$ ein:
\[
    S_0\to S.
\]

\textbf{Schritt 2: $\varepsilon$-Regeln entfernen.}

Nullable sind $A$, $B$ und $S$, denn
\[
    A\Rightarrow\varepsilon,\qquad B\Rightarrow A\Rightarrow\varepsilon,\qquad
    S\Rightarrow A\Rightarrow\varepsilon.
\]
Aus $S\to SaSb$ entstehen daher alle Varianten, in denen die beiden Vorkommen von $S$
optional gestrichen werden. Nur fuer die neue Startvariable behalten wir $\varepsilon$:
\[
\begin{aligned}
S_0 &\to S \mid \varepsilon,\\
S &\to SaSb \mid aSb \mid Sab \mid ab \mid A,\\
A &\to B \mid a,\\
B &\to bB \mid b \mid A \mid C,\\
C &\to c.
\end{aligned}
\]

\textbf{Schritt 3: Kettenregeln entfernen.}

Die Kettenregeln sind $S_0\to S$, $S\to A$, $A\to B$, $B\to A$ und $B\to C$.
Wir ersetzen sie, indem wir die nichttrivialen Regeln der erreichbaren Variablen direkt uebernehmen:
\[
\begin{aligned}
S_0 &\to \varepsilon \mid SaSb \mid aSb \mid Sab \mid ab
        \mid a \mid bB \mid b \mid c,\\
S &\to SaSb \mid aSb \mid Sab \mid ab
        \mid a \mid bB \mid b \mid c,\\
A &\to a \mid bB \mid b \mid c,\\
B &\to bB \mid a \mid b \mid c,\\
C &\to c.
\end{aligned}
\]
Die Variablen $A$ und $C$ sind nun vom Startsymbol aus nicht mehr erreichbar und koennen entfernt
werden.

\textbf{Schritt 4: Terminale in langen Regeln ersetzen und binaer machen.}

Wir fuehren
\[
    T_a\to a,\qquad T_b\to b
\]
ein. Fuer die langen rechten Seiten nutzen wir zusaetzlich
\[
    E\to T_aT_b,\qquad G\to ST_b,\qquad F\to T_aG.
\]
Damit erhalten wir die CNF-Grammatik
\[
G_{\mathrm{CNF}}=(V',\{a,b,c\},R',S_0)
\]
mit
\[
V'=\{S_0,S,B,T_a,T_b,E,F,G\}
\]
und Regeln
\[
\begin{aligned}
S_0 &\to \varepsilon \mid SF \mid T_aG \mid SE \mid T_aT_b
        \mid a \mid T_bB \mid b \mid c,\\
S &\to SF \mid T_aG \mid SE \mid T_aT_b
        \mid a \mid T_bB \mid b \mid c,\\
B &\to T_bB \mid a \mid b \mid c,\\
E &\to T_aT_b,\\
F &\to T_aG,\\
G &\to ST_b,\\
T_a &\to a,\\
T_b &\to b.
\end{aligned}
\]
Alle rechten Seiten haben nun die Form $XY$, ein einzelnes Terminal oder, nur bei $S_0$, das Wort
$\varepsilon$. Ausserdem kommt $S_0$ auf keiner rechten Seite vor.

\section*{Aufgabe 5.7}

\taskbox{
Zeigen Sie, dass die kontextfreien Sprachen unter Vereinigung, Konkatenation und Kleene-Stern
abgeschlossen sind. Geben Sie jeweils eine Konstruktionsvorschrift fuer die entsprechende Grammatik an.
}

\methodbox{
Eine Sprache ist kontextfrei, wenn es eine CFG gibt, die sie erzeugt. Fuer Abschlussbeweise reicht es
daher, aus gegebenen CFGs neue CFGs zu konstruieren, die genau die gewuenschte zusammengesetzte
Sprache erzeugen.
}

Seien $L_1,L_2$ kontextfrei. Dann gibt es CFGs
\[
G_1=(V_1,\Sigma,R_1,S_1),\qquad
G_2=(V_2,\Sigma,R_2,S_2)
\]
mit $L(G_1)=L_1$ und $L(G_2)=L_2$. Wir nehmen ohne Einschraenkung an, dass
$V_1\cap V_2=\emptyset$ ist, denn Variablen kann man umbenennen.

\textbf{Vereinigung.}
Definiere
\[
G_\cup=(V_1\cup V_2\cup\{S\},\Sigma,R_1\cup R_2\cup\{S\to S_1\mid S_2\},S).
\]
Dann waehlt die erste Regel genau aus, ob ein Wort mit $G_1$ oder mit $G_2$ erzeugt wird. Also gilt
\[
L(G_\cup)=L_1\cup L_2.
\]

\textbf{Konkatenation.}
Definiere
\[
G_\circ=(V_1\cup V_2\cup\{S\},\Sigma,R_1\cup R_2\cup\{S\to S_1S_2\},S).
\]
Jede Ableitung erzeugt zuerst ein Wort aus $L_1$ und danach ein Wort aus $L_2$. Umgekehrt kann jedes
$uv$ mit $u\in L_1$ und $v\in L_2$ so abgeleitet werden. Also gilt
\[
L(G_\circ)=L_1L_2.
\]

\textbf{Kleene-Stern.}
Fuer $L_1^*$ definieren wir
\[
G_*=(V_1\cup\{S\},\Sigma,R_1\cup\{S\to S_1S\mid \varepsilon\},S).
\]
Die Regel $S\to\varepsilon$ erzeugt null Faktoren. Jede Anwendung von $S\to S_1S$ fuegt vorne ein
weiteres Wort aus $L_1$ hinzu. Daher erzeugt die Grammatik genau endliche Konkatenationen von
Woertern aus $L_1$, also
\[
L(G_*)=L_1^*.
\]

\section*{Aufgabe 5.8}

\taskbox{
Bestimmen Sie die Sprache, die von dem gegebenen PDA akzeptiert wird, und geben Sie eine moeglichst
intuitive Beschreibung an.
}

\methodbox{
Bei einem PDA bedeutet die Kantenbeschriftung $a,b\to c$: Lies $a$, entferne $b$ vom Stack und lege
$c$ auf den Stack. Falls ein Eintrag $\varepsilon$ ist, wird an dieser Stelle nichts gelesen, entfernt oder
hinzugefuegt. Ein nichtdeterministischer PDA akzeptiert, wenn mindestens ein Berechnungspfad in einem
akzeptierenden Zustand endet, nachdem die Eingabe gelesen wurde.
}

Der PDA legt zuerst einen Bodenmarker $\$$ auf den Stack. Danach gilt im Zustand $q_1$:
\[
\texttt{(},\varepsilon\to 1
\qquad\text{und}\qquad
\texttt{)},1\to\varepsilon.
\]
Also wird fuer jede oeffnende Klammer ein $1$ auf den Stack gelegt, und fuer jede schliessende Klammer
wird ein $1$ entfernt. Akzeptiert wird genau dann, wenn am Ende wieder nur der Bodenmarker uebrig ist.

Damit ist die akzeptierte Sprache
\[
\boxed{L(P)=D}
\]
mit
\[
\begin{aligned}
D=\{w\in\{\texttt{(},\texttt{)}\}^* \mid{}&
\text{jeder Praefix von }w\text{ hat mindestens so viele \texttt{(} wie \texttt{)},}\\
&\text{und insgesamt sind es gleich viele}\}.
\end{aligned}
\]
Intuitiv ist das die Sprache der korrekt geklammerten Woerter mit genau einer Klammerart.

\examplebox{
Akzeptiert werden zum Beispiel
\[
\varepsilon,\qquad \texttt{()},\qquad \texttt{(())},\qquad \texttt{()()}.
\]
Nicht akzeptiert werden zum Beispiel
\[
\texttt{)(},\qquad \texttt{(()},\qquad \texttt{())}.
\]
}

\section*{Aufgabe 5.9}

\taskbox{
Entwerfen Sie einen PDA fuer
\[
L=\{a^ib^jc^k\mid 2j=3k \ \lor\ i+k=j\}
\]
und begruenden Sie kurz, weshalb der PDA die Sprache erkennt.
}

\methodbox{
Da PDAs nichtdeterministisch sind, kann man fuer eine Vereinigung zweier Bedingungen zuerst
nichtdeterministisch einen Zweig waehlen. Der eine Zweig prueft $2j=3k$, der andere Zweig prueft
$i+k=j$.
}

Wir beschreiben einen PDA $P$ ueber dem Eingabealphabet $\{a,b,c\}$. Zuerst legt $P$ den Bodenmarker
$\$$ auf den Stack und waehlt dann nichtdeterministisch einen der beiden folgenden Zweige.

\subsection*{Zweig 1: Pruefe $2j=3k$}

Die $a$'s werden ignoriert. Fuer jedes gelesene $b$ werden zwei Symbole $X$ auf den Stack gelegt. Fuer
jedes gelesene $c$ werden drei Symbole $X$ entfernt. Am Ende wird akzeptiert, wenn genau der
Bodenmarker erreicht ist.

Eine moegliche Uebergangsbeschreibung ist:
\[
\begin{array}{rcl}
r_0 &:& a,\varepsilon\to\varepsilon \text{ bleibt in } r_0,\\
r_0 &:& \varepsilon,\varepsilon\to\varepsilon \text{ nach } r_b,\\
r_b &:& b,\varepsilon\to X \text{ nach } r_{b,2},\\
r_{b,2} &:& \varepsilon,\varepsilon\to X \text{ nach } r_b,\\
r_b &:& c,X\to\varepsilon \text{ nach } r_{c,1},\\
r_{c,1} &:& \varepsilon,X\to\varepsilon \text{ nach } r_{c,2},\\
r_{c,2} &:& \varepsilon,X\to\varepsilon \text{ nach } r_c,\\
r_c &:& c,X\to\varepsilon \text{ nach } r_{c,1},\\
r_b,r_c &:& \varepsilon,\$\to\varepsilon \text{ nach } q_{\mathrm{acc}}.
\end{array}
\]
Dieser Zweig akzeptiert genau dann, wenn die Anzahl der gelegten $X$ gleich der Anzahl der entfernten
$X$ ist, also genau dann, wenn $2j=3k$ gilt.

\subsection*{Zweig 2: Pruefe $i+k=j$}

Fuer jedes $a$ wird ein Symbol $A$ auf den Stack gelegt. In der $b$-Phase werden zuerst diese
$A$-Symbole entfernt. Sobald keine $A$-Symbole mehr uebrig sind, wird fuer jedes weitere $b$ ein
Symbol $B$ auf den Stack gelegt. In der $c$-Phase wird fuer jedes $c$ ein $B$ entfernt. Am Ende wird
akzeptiert, wenn genau der Bodenmarker erreicht ist.

Eine moegliche Uebergangsbeschreibung ist:
\[
\begin{array}{rcl}
s_a &:& a,\varepsilon\to A \text{ bleibt in } s_a,\\
s_a &:& b,A\to\varepsilon \text{ nach } s_{\mathrm{pop}},\\
s_{\mathrm{pop}} &:& b,A\to\varepsilon \text{ bleibt in } s_{\mathrm{pop}},\\
s_a,s_{\mathrm{pop}} &:& \varepsilon,\$\to \$ \text{ nach } s_{\mathrm{push}},\\
s_{\mathrm{push}} &:& b,\varepsilon\to B \text{ bleibt in } s_{\mathrm{push}},\\
s_{\mathrm{push}} &:& c,B\to\varepsilon \text{ nach } s_c,\\
s_c &:& c,B\to\varepsilon \text{ bleibt in } s_c,\\
s_{\mathrm{push}},s_c &:& \varepsilon,\$\to\varepsilon \text{ nach } q_{\mathrm{acc}}.
\end{array}
\]
Dieser Zweig akzeptiert genau dann, wenn die Zahl der $b$'s ausreicht, um zuerst alle $i$ Symbole fuer
die $a$'s zu entfernen, und danach genau $k$ weitere $b$'s fuer die $c$'s gespeichert wurden. Also gilt
genau
\[
    j=i+k.
\]

Da der PDA zu Beginn nichtdeterministisch einen der beiden Zweige waehlt, akzeptiert er genau dann,
wenn mindestens eine der Bedingungen erfuellt ist. Somit gilt
\[
\boxed{L(P)=\{a^ib^jc^k\mid 2j=3k \ \lor\ i+k=j\}.}
\]

\end{document}
